\documentclass[10pt, a4paper]{article}

\usepackage[left=1.15cm, right=1.15cm, top=1.1cm, bottom=1.1cm]{geometry}
\usepackage{titlesec}
\usepackage{enumitem}
\usepackage{xcolor}
\usepackage{hyperref}
\usepackage{tabularx}

\definecolor{citeColor}{RGB}{0,20,115}
\hypersetup{
    colorlinks=true,
    linkcolor=citeColor,
    urlcolor=citeColor,
    citecolor=citeColor
}

\setlength{\parindent}{0pt}
\setlength{\tabcolsep}{0pt}
\renewcommand{\baselinestretch}{1.03}
\raggedright
\pagestyle{empty}

\titleformat{\section}
  {\vspace{-1pt}\normalsize\scshape\raggedright}
  {}
  {0em}
  {}
  [\color{black}\titlerule]

\titlespacing*{\section}{0pt}{0.85em}{0.35em}
\setlist[itemize]{leftmargin=1.35em, nosep, itemsep=0.2em, topsep=0.15em}
\setlist[enumerate]{leftmargin=1.6em, nosep, itemsep=0.28em, topsep=0.15em}

\newcommand{\resumeSubheading}[4]{
  \vspace{0.15em}
  \begin{tabular*}{\textwidth}{l@{\extracolsep{\fill}}r}
    \textbf{#1} & #2 \\
    \textit{\small #3} & \textit{\small #4} \\
  \end{tabular*}
  \vspace{-0.35em}
}

\newcommand{\workHeading}[3]{
  \vspace{0.25em}
  \begin{tabular*}{\textwidth}{l@{\extracolsep{\fill}}r}
    \textbf{#1} & \small #2 \\
    \textit{\small #3} & \\
  \end{tabular*}
  \vspace{-0.45em}
}

\newcommand{\projectEntry}[3]{
  \vspace{0.28em}
  \begin{tabularx}{\textwidth}{@{}>{\raggedright\arraybackslash}X r@{}}
    \textbf{#1} & \small #3 \\
    \textit{\small #2} & \\
  \end{tabularx}
  \vspace{-0.25em}
}

\begin{document}

{\LARGE \textbf{Jiaxuan Wang}}\\[0.3em]
\textbf{Email:} \href{mailto:wangjx_hust@hust.edu.cn}{wangjx\_hust@hust.edu.cn}\\
\textbf{Phone:} +86 188 7996 5616\\
\href{https://github.com/xuanhh567}{[GitHub: xuanhh567]}

\section{Short Bio}
I am currently an undergraduate student in Electronic Information Engineering at Huazhong University of Science and Technology, where I serve as Captain of Dian Team and Lead of its Robotics Group.

I closely follow frontier research in embodied AI and robot learning, with a strong interest in identifying practical problems and formulating feasible technical ideas. I combine research intuition with hands-on engineering by reproducing representative models, adapting them to robotic training workflows, and validating ideas on open-source robotic arms. My long-term goal is to build general-purpose robots that learn robust skills and assist people in the real world.

\section{Research Interests}
Embodied AI; Robot Learning; Real-World Robot Deployment.

\section{Education}
\begin{tabularx}{\textwidth}{@{}>{\raggedright\arraybackslash}X r@{}}
\textbf{Huazhong University of Science and Technology} & Sep. 2023 -- Jun. 2027 \\
B.Eng. in Electronic Information Engineering & \\
\end{tabularx}
Project-Based Information Engineering Experimental Class (``Seed Class''); Weighted Average: \textbf{84.5/100}; CET-6: \textbf{519}.

\section{Competitions and Projects}
\projectEntry
  {ICRA 2025 Sim2Real Challenge}
  {Core Member, DianRobot}
  {2025}
\begin{itemize}
    \item Designed and evaluated tabletop bimanual grasping strategies for objects of different shapes and categories, including action-sequence debugging, rollout analysis, real-robot validation, and challenge demonstration preparation.
\end{itemize}
\projectEntry
  {Robotic Arm Application for Chemistry Experiments Based on Imitation Learning}
  {Provincial Undergraduate Innovation Project; Team Leader}
  {2025}
\begin{itemize}
    \item Led task decomposition, data-collection workflow design, imitation-learning training, real-arm testing, and grasping experiments for chemistry-lab instruments.
\end{itemize}
\projectEntry
  {Robot Coffee-Making Based on Reinforcement Learning}
  {Industry Collaboration Project with CloudMinds Robotics}
  {Jul. 2023 -- Apr. 2024}
\begin{itemize}
    \item Built the coffee-service demonstration workflow; supported cup grasping, localization, and success-check modules; addressed joint zero-drift through visual compensation.
\end{itemize}

\section{Publications}
\begin{enumerate}
    \item \textbf{[ICTC 2026]} Yefan Lin, \textbf{Jiaxuan Wang}, Baoyu Liu, Zheng Lv, and Xiaojun Hei, ``Pallet-ACT: A Unified Data-Training-Inference Framework for Industrial Robotic Palletizing.'' \textit{7th Information Communication Technologies Conference}, Nanjing, China, May 2026. \textbf{Accepted}.
    \item \textbf{[ROBIO 2025]} Yuhan Huang, Zijian Ning, Ziyuan Wang, \textbf{Jiaxuan Wang}, Xiaojun Hei, ``Reinforcement Learning with Stage-Wise Model Fusion for Sequential Object Manipulation Tasks.'' \textit{IEEE International Conference on Robotics and Biomimetics}, Chengdu, China, Dec. 2025.
\end{enumerate}

\section{Awards}
\begin{itemize}
    \item \textbf{ICRA 2025 Sim2Real Challenge:} International Second Prize, 4th Place Worldwide.
    \item \textbf{Chinese Collegiate Computing Competition (4C), Central-South Division, Undergraduate Group:} Second Prize, 2024, for ``Reinforcement-Learning-Enabled Robotic Coffee Service Application.''
    \item \textbf{Robotic Arm Application for Chemistry Experiments Based on Imitation Learning:} Provincial Undergraduate Innovation Project, Team Leader, rated Excellent, 2025.
    \item \textbf{Dian Team Reception Robot Application Design:} University-level Undergraduate Innovation Project, rated Excellent, 2023.
\end{itemize}

\section{Technical Skills}
\begin{itemize}
    \item \textbf{Robot Learning:} reinforcement learning, imitation learning, robotic control, Pi-series vision-language-action (VLA) model reproduction.
    \item \textbf{Robotics Stack:} 6-axis robotic-arm operation and deployment, ROS 2, MoveIt 2, MuJoCo simulation, real-robot data collection.
    \item \textbf{Engineering:} Python, C, PyTorch, NumPy, OpenCV, Git, Linux, SolidWorks, 3D printing.
\end{itemize}

\end{document}
